\documentclass[12pt]{article}

\usepackage{amsmath}
\usepackage{amssymb}
\usepackage{geometry}
\usepackage{setspace}
\usepackage{titlesec}

\geometry{margin=1in}
\setstretch{1.2}

\title{CSE400 -- Fundamentals of Probability in Computing\\
Lecture Scribe: Joint Random Variables and Joint Distributions (L15--L18)}
\author{Instructor: Dhaval Patel, PhD}
\date{}

\begin{document}

\maketitle

\section*{Course Information}

\textbf{Course:} CSE400 -- Fundamentals of Probability in Computing \\
\textbf{Topic:} Joint Random Variables, Joint PMF, Joint PDF, and Joint CDF

\hrule
\vspace{0.5cm}

\section{Motivation: Why Study Joint Random Variables?}

In many scientific and engineering problems, we are interested in \textbf{more than one random variable at the same time}. The value of one variable may depend on or be related to the value of another.

Thus, instead of studying a single random variable independently, we study \textbf{two random variables together} using \textbf{joint probability distributions}.

The purpose of studying joint random variables is therefore:

\begin{itemize}
\item To analyze how two random variables behave \textbf{simultaneously}
\item To determine whether there exists a \textbf{relationship or correlation}
\item To compute probabilities involving \textbf{combined events}
\end{itemize}

\section{Motivating Examples}

\subsection{Smart Transportation Example}

Consider a smart traffic system with two random variables:

\begin{itemize}
\item $X =$ Vehicle speed
\item $Y =$ Traffic density
\end{itemize}

Key question:

\begin{quote}
Can vehicle speed be high when traffic density is also high?
\end{quote}

Possible answers: \textbf{Yes / No}

This indicates that the values of $X$ and $Y$ may be \textbf{correlated}.

Traffic conditions affect vehicle speed. Therefore, analyzing only one variable does not provide complete information.

Applications include:

\begin{itemize}
\item Adaptive traffic signals
\item Congestion prediction
\item Autonomous vehicle routing
\end{itemize}

\subsection{Wireless Network Example}

Consider a wireless communication system.

Define:

\begin{itemize}
\item $X =$ Signal-to-Noise Ratio (SNR)
\item $Y =$ Data throughput
\end{itemize}

Question:

\begin{quote}
If SNR is high, will throughput always be high?
\end{quote}

The answer is \textbf{not necessarily}.

Even if SNR is high, \textbf{network congestion} may reduce throughput.

Applications include:

\begin{itemize}
\item 5G scheduling
\item Adaptive modulation
\item Network planning
\end{itemize}

\subsection{Additional Examples}

\textbf{Weather Prediction}

\begin{itemize}
\item $X =$ Rainfall
\item $Y =$ Temperature
\end{itemize}

Applications: Weather forecasting.

\vspace{0.3cm}

\textbf{Drone Delivery Systems}

\begin{itemize}
\item $X =$ Wind speed
\item $Y =$ Delivery time
\end{itemize}

Applications: Logistics and UAV systems.

\vspace{0.3cm}

\textbf{Rolling Two Dice}

\begin{itemize}
\item $X =$ outcome of the first die
\item $Y =$ outcome of the second die
\end{itemize}

Example questions:

\begin{itemize}
\item $X + Y = 7$
\item $X > Y$
\end{itemize}

These require the \textbf{joint distribution}.

\section{Discrete Case -- Joint Probability Mass Function (Joint PMF)}

When random variables are \textbf{discrete}, their joint behavior is described by the \textbf{Joint Probability Mass Function}.

The joint PMF is written as

\[
P(X = x, Y = y)
\]

Interpretation:

\begin{itemize}
\item $X$ takes value $x$
\item $Y$ simultaneously takes value $y$
\end{itemize}

Thus, the joint PMF gives the probability of \textbf{simultaneous outcomes of two random variables}.

\section{Joint PMF Table Representation}

For discrete random variables, the joint PMF is commonly represented using a \textbf{table}.

\begin{itemize}
\item Each \textbf{row} corresponds to a value of $X$
\item Each \textbf{column} corresponds to a value of $Y$
\item Each cell contains $P(X=x,Y=y)$
\end{itemize}

This table provides a \textbf{complete description of the joint distribution}.

\section{Example: Rolling Two Dice}

Consider the experiment of rolling \textbf{two fair dice}.

Define:

\begin{itemize}
\item $X =$ value on the first die
\item $Y =$ value on the second die
\end{itemize}

Possible values:

\[
X,Y \in \{1,2,3,4,5,6\}
\]

Total outcomes:

\[
6 \times 6 = 36
\]

Since both dice are fair:

\[
P(X=x,Y=y)=\frac{1}{36}
\]

for all valid pairs.

Examples:

\begin{itemize}
\item Probability that $X+Y=7$
\item Probability that $X>Y$
\end{itemize}

These require multiple joint outcomes.

\section{Continuous Case -- Joint Probability Density Function (Joint PDF)}

If the random variables are \textbf{continuous}, the joint distribution is represented by a \textbf{Joint Probability Density Function (PDF)}.

Notation:

\[
f_{X,Y}(x,y)
\]

This represents the \textbf{probability density} at the point $(x,y)$.

Probabilities are computed by integration.

For a region $R$:

\[
P((X,Y)\in R) = \iint_R f_{X,Y}(x,y)\,dx\,dy
\]

Thus, probability corresponds to the \textbf{integral of the density over the region}.

\section{Geometric Interpretation of Joint PDF}

The function $f_{X,Y}(x,y)$ forms a \textbf{surface over the $(x,y)$ plane}.

\begin{itemize}
\item The plane represents possible values of $X$ and $Y$
\item Probability corresponds to the \textbf{volume under the surface}
\end{itemize}

\section{Joint Cumulative Distribution Function (Joint CDF)}

Definition:

\[
F_{X,Y}(x,y) = P(X \le x,\; Y \le y)
\]

Interpretation:

\begin{itemize}
\item $X$ is less than or equal to $x$
\item $Y$ is less than or equal to $y$
\end{itemize}

It accumulates probability in the \textbf{lower-left region} of the coordinate plane.

\section{Joint CDF -- Continuous Case}

For continuous random variables:

\[
F_{X,Y}(x,y) =
\int_{-\infty}^{x}
\int_{-\infty}^{y}
f_{X,Y}(u,v)\,du\,dv
\]

Reasoning:

\begin{enumerate}
\item Consider all $u \le x$
\item Consider all $v \le y$
\item Integrate the joint density over the region
\end{enumerate}

\section{Joint CDF -- Discrete Case}

For discrete variables:

\[
F_{X,Y}(x,y) =
\sum_{u\le x}
\sum_{v\le y}
P(X=u,Y=v)
\]

Steps:

\begin{enumerate}
\item Identify values $u \le x$
\item Identify values $v \le y$
\item Add probabilities for all such pairs
\end{enumerate}

\section{Properties of the Joint CDF}

\subsection{Property 1: Probability Range}

\[
0 \le F_{X,Y}(x,y) \le 1
\]

\subsection{Property 2: Monotonicity}

The joint CDF is \textbf{non-decreasing} in both $x$ and $y$.

Increasing either variable enlarges the region where probability is accumulated.

\subsection{Property 3: Boundary Behavior}

As $x$ and $y$ increase without bound:

\[
F_{X,Y}(x,y) \rightarrow 1
\]

When $x$ or $y \rightarrow -\infty$:

\[
F_{X,Y}(x,y) \rightarrow 0
\]

\section{Conclusion}

The lecture introduces the concept of \textbf{joint random variables}, which allow the probabilistic analysis of systems involving multiple uncertain quantities.

Key concepts covered:

\begin{itemize}
\item Motivation for studying joint random variables
\item Joint PMF for discrete variables
\item Joint PDF for continuous variables
\item Geometric interpretation of joint densities
\item Definition and computation of the Joint CDF
\item Properties of the Joint CDF
\end{itemize}

These concepts form the mathematical foundation for analyzing probabilistic relationships between multiple variables in real-world systems such as transportation networks, wireless communication systems, weather prediction models, and stochastic experiments.

\end{document}